% !TeX root = ../../Book3.tex
\section{\texorpdfstring{\(\theta\)}{\AnalysisPdfTheta} 级数作为生成函数}
\label{b3:sec:2403}

系数计数必须先约定是否区分顺序和符号。theta 乘法自动区分各坐标，而且正负整数分别出现，因此它计算的是有序带符号表示数。

\subsection{生成函数}
\label{b3:subsec:240301}

设 \(r_k(n)=\#\{(a_1,\ldots,a_k)\in\mathbb Z^k:\sum a_j^2=n\}\)。每个坐标绝对值不超过 \(\sqrt n\)，故数目有限。对 \(|q|<1\)，绝对收敛保证
\[
\theta_3(0,\tau)^k=\sum_{n\ge0}r_k(n)q^n.
\]
此处 \(r_k(0)=1\)。例如 \(r_2(1)=4\)、\(r_4(1)=8\)，分别选择非零坐标的位置和它的符号，不能把它们解释为不计顺序的分拆数。

\subsection{乘积与卷积}
\label{b3:subsec:240302}

乘法给 \(r_{k+\ell}(n)=\sum_{m=0}^n r_k(m)r_\ell(n-m)\)。这是把前 \(k\) 个坐标与后 \(\ell\) 个坐标分别计数的精确分解。

\ProseParagraphBreak

若改用 \(\theta_4\)，每个坐标带符号 \((-1)^{a_j}\)。由于 \(a_j^2\equiv a_j\pmod2\)，同一层 \(\sum a_j^2=n\) 上总符号恒为 \((-1)^n\)。因此 \(\theta_4^k=\sum(-1)^nr_k(n)q^n\)。它筛选奇偶层，却不单独筛选每个坐标的奇偶模式。

\subsection{系数比较}
\label{b3:subsec:240303}

在四个变量中还有一个有用的精确变换。对 \(a=(a_1,a_2,a_3,a_4)\) 用正交矩阵
\[
\frac12\begin{pmatrix}1&1&1&1\\1&1&-1&-1\\1&-1&1&-1\\1&-1&-1&1\end{pmatrix}
\]
变换为 \(b\)。当 \(\sum a_j\) 为奇数时，所有 \(b_j\in\mathbb Z+1/2\)；反向变换相同。半整数四元组又分为两类：写 \(b_j=m_j+1/2\)，按 \(\sum m_j\) 的奇偶分。只有其中一类在逆变换下给整数向量；将一个 \(b_j\) 改号交换两类且保持平方和。故全部半整数向量的计数恰为奇数平方和整数向量计数的两倍。比较级数得到
\begin{equation}\label{b3:eq:24-quartic-identity}
\theta_3^4=\theta_4^4+\theta_2^4.
\end{equation}
这个推导既说明常数二从何而来，也说明为什么四次方出现。它不是把前几项相同当成无穷恒等式，而是给出保持范数的对应。

\ProseParagraphBreak

生成函数还可以筛选同余条件。若 \(A(q)=\sum a_nq^n\)，\(\zeta_m=e^{2\pi\mathrm{i}/m}\)，则有限 Fourier 正交关系给
\[
 \frac1m\sum_{j=0}^{m-1}\zeta_m^{-bj}A(\zeta_m^j q)
 =\sum_{n\equiv b\ (m)}a_nq^n.
\]
交换有限和与级数不需要新的收敛定理；只要 \(|q|<1\)，每个旋转参数仍在原收敛圆盘内。对平方和而言，这个操作筛选总平方和的剩余类。若要每个坐标分别满足某些同余条件，则须先对每个一维 theta 因子筛选，再相乘。两个程序统计的集合不同，不能仅凭最后都含若干符号因子就互换。
\ProseParagraphBreak
四维正交变换中，半整数层的两类也可直接检查。写 \(b_j=m_j+1/2\)，逆变换的第一坐标为 \((\sum m_j+2)/2\)，其余坐标具有同样的整数奇偶条件。因此恰在 \(\sum m_j\) 偶数时，逆像为整数四元组；这时逆像的坐标和为 \(2b_1\)，必为奇数。把一个 \(b_j\) 改号，相应的整数 \(m_j\) 变成 \(-m_j-1\)，总和奇偶翻转，所以两类等势且逐层保范数。这样便把前述系数比较中的二倍关系落实为显式双射，而非只比较总的无穷集合大小。
\ProseParagraphBreak
平均表示数可以从几何体积作一个独立的数量级检验。令 \(A_k(R)\) 为 \(\mathbb Z^k\) 中长度不超过 \(R\) 的点数。以每个格点为中心放置边长一的闭立方体，它们内部不交，每个体积一。所取立方体的并包含半径 \(R-\sqrt{k}/2\) 的球，又包含在半径 \(R+\sqrt{k}/2\) 的球内。因此
\[
 A_k(R)=V_kR^k+O_k(R^{k-1}),
\]
其中 \(V_k\) 为单位球体积。这个估计不控制单个球壳上的点数，却控制从0累计到 \(R^2\) 的表示数。令 \(X=R^2\)，即为 \(\sum_{n\le X}r_k(n)=V_kX^{k/2}+O_k(X^{(k-1)/2})\)。平方和问题中，逐项精确公式与累计体积估计是不同层次的结果。
